% GENERATED FILE. Edit canonical lessons, then run npm run generate:books.
% canonical-source-hash: fnv1a64:c8b16b9f3a2e1d2d
% canonical-lessons: TE-C516-cupu, TE-C516-vinikidi, TE-C516-cevudu, TE-C516-alerji, TE-C516-nidralemi

\chapter{Senses}
\label{ch:te-senses}
\hlchaptermodality{516}

\begin{chapteropening}
\textbf{By the end of this chapter:} \emph{I can talk about eyesight, hearing, allergies and sleeplessness.}

Complete the last lesson of chapter 516: I can talk about eyesight, hearing, allergies and sleeplessness.
\end{chapteropening}

\section[\texorpdfstring{\te{చూపు} (cūpu)}{చూపు (cūpu)}]{\te{చూపు} (cūpu) — eyesight; a look}

\label{lesson:TE-C516-cupu}



\begin{quote}
Before the new one: say the Telugu for the spine, then the Telugu for breath.
\end{quote}

\subsection*{You'll want to know: \te{చూపు}}
\textbf{\te{చూపు}} — \emph{cūpu} — \textquotedblleft{}eyesight; a look\textquotedblright{}.

\begin{quote}
\te{వయసుతో} \te{తాతయ్య} \te{చూపు} \te{తగ్గింది}.
\end{quote}

\begin{quote}
\emph{vayasutō tātayya cūpu taggindi.}
\end{quote}

\begin{quote}
With age grandpa's eyesight has weakened.
\end{quote}

Also a glance: \te{ఒక} \te{చూపు}, one look. The noun of \te{చూడు}, to see.

\begin{cousinweb}
\textbf{\te{చూపు}} (\emph{cūpu}) is native Dravidian: DEDR 2735 lists Telugu \emph{cūcu}, \textquotedblleft{}to see\textquotedblright{}, and \emph{cūpu}, \textquotedblleft{}sight, vision, look, glance\textquotedblright{}.
\end{cousinweb}

\begin{grammarlens}[title={What you've built}]
One more word for describing what you lived and felt.
\end{grammarlens}

\subsection*{Your turn}
\begin{itemize}
\raggedright
  \item \emph{Say it:} \emph{cūpu}
  \item \emph{Say it:} \emph{cūpu}, once more
  \item \emph{Say it:} \emph{vayasutō tātayya cūpu taggindi.} — With age grandpa's eyesight has weakened.
  \item \emph{Recall:} say the Telugu for the spine, then the Telugu for breath, then say \emph{cūpu} again
\end{itemize}

\begin{tcolorbox}[breakable,colback=teal!4,colframe=teal!35!black,title={Before you move on}]
What does \te{చూపు} mean? (\textquotedblleft{}Eyesight; a look\textquotedblright{}.) Say it once more.
\end{tcolorbox}
\section[\texorpdfstring{\te{వినికిడి} (vinikiḍi)}{వినికిడి (vinikiḍi)}]{\te{వినికిడి} (vinikiḍi) — hearing}

\label{lesson:TE-C516-vinikidi}



\begin{quote}
Before the new one: say the Telugu for breath, then the Telugu for eyesight.
\end{quote}

\subsection*{You'll want to know: \te{వినికిడి}}
\textbf{\te{వినికిడి}} — \emph{vinikiḍi} — \textquotedblleft{}hearing\textquotedblright{}.

\begin{quote}
\te{పెద్ద} \te{శబ్దాల} \te{వల్ల} \te{వినికిడి} \te{దెబ్బతింటుంది}.
\end{quote}

\begin{quote}
\emph{pedda śabdāla valla vinikiḍi debbatiṇṭundi.}
\end{quote}

\begin{quote}
Loud noises damage the hearing.
\end{quote}

The noun of \te{విను}, to hear.

\begin{cousinweb}
\textbf{\te{వినికిడి}} (\emph{vinikiḍi}) is native Dravidian: DEDR 5516 lists Telugu \emph{vinu}, \textquotedblleft{}to hear, to listen\textquotedblright{}, and \emph{vinikiḍi}, \textquotedblleft{}hearing, listening\textquotedblright{}.
\end{cousinweb}

\begin{grammarlens}[title={What you've built}]
One more word for describing what you lived and felt.
\end{grammarlens}

\subsection*{Your turn}
\begin{itemize}
\raggedright
  \item \emph{Say it:} \emph{vinikiḍi}
  \item \emph{Say it:} \emph{vinikiḍi}, once more
  \item \emph{Say it:} \emph{pedda śabdāla valla vinikiḍi debbatiṇṭundi.} — Loud noises damage the hearing.
  \item \emph{Recall:} say the Telugu for breath, then the Telugu for eyesight, then say \emph{vinikiḍi} again
\end{itemize}

\begin{tcolorbox}[breakable,colback=teal!4,colframe=teal!35!black,title={Before you move on}]
What does \te{వినికిడి} mean? (\textquotedblleft{}Hearing\textquotedblright{}.) Say it once more.
\end{tcolorbox}
\section[\texorpdfstring{\te{చెవుడు} (cevuḍu)}{చెవుడు (cevuḍu)}]{\te{చెవుడు} (cevuḍu) — deafness}

\label{lesson:TE-C516-cevudu}



\begin{quote}
Before the new one: say the Telugu for eyesight, then the Telugu for hearing.
\end{quote}

\subsection*{You'll want to know: \te{చెవుడు}}
\textbf{\te{చెవుడు}} — \emph{cevuḍu} — \textquotedblleft{}deafness\textquotedblright{}.

\begin{quote}
\te{ముసలితనంలో} \te{ఆయనకు} \te{కొంచెం} \te{చెవుడు} \te{వచ్చింది}.
\end{quote}

\begin{quote}
\emph{musalitanaṁlō āyanaku koñceṁ cevuḍu vaccindi.}
\end{quote}

\begin{quote}
In old age he became a little deaf.
\end{quote}

Built on \te{చెవి}, ear. \te{చెవిటివాడు}, a deaf man.

\begin{cousinweb}
\textbf{\te{చెవుడు}} (\emph{cevuḍu}) is built on the native \te{చెవి}, \textquotedblleft{}ear\textquotedblright{} (DEDR 1977, with Tamil \emph{cevi}), with a noun ending; DEDR does not list \emph{cevuḍu} itself.
\end{cousinweb}

\begin{grammarlens}[title={What you've built}]
One more word for describing what you lived and felt.
\end{grammarlens}

\subsection*{Your turn}
\begin{itemize}
\raggedright
  \item \emph{Say it:} \emph{cevuḍu}
  \item \emph{Say it:} \emph{cevuḍu}, once more
  \item \emph{Say it:} \emph{musalitanaṁlō āyanaku koñceṁ cevuḍu vaccindi.} — In old age he became a little deaf.
  \item \emph{Recall:} say the Telugu for eyesight, then the Telugu for hearing, then say \emph{cevuḍu} again
\end{itemize}

\begin{tcolorbox}[breakable,colback=teal!4,colframe=teal!35!black,title={Before you move on}]
What does \te{చెవుడు} mean? (\textquotedblleft{}Deafness\textquotedblright{}.) Say it once more.
\end{tcolorbox}
\section[\texorpdfstring{\te{అలెర్జీ} (alerjī)}{అలెర్జీ (alerjī)}]{\te{అలెర్జీ} (alerjī) — an allergy}

\label{lesson:TE-C516-alerji}



\begin{quote}
Before the new one: say the Telugu for hearing, then the Telugu for deafness.
\end{quote}

\subsection*{You'll want to know: \te{అలెర్జీ}}
\textbf{\te{అలెర్జీ}} — \emph{alerjī} — \textquotedblleft{}an allergy\textquotedblright{}.

\begin{quote}
\te{నాకు} \te{దుమ్ము} \te{అంటే} \te{అలెర్జీ}.
\end{quote}

\begin{quote}
\emph{nāku dummu aṇṭē alerjī.}
\end{quote}

\begin{quote}
I am allergic to dust.
\end{quote}

-\te{అంటే} \te{అలెర్జీ}, allergic to.

\begin{cousinweb}
\textbf{\te{అలెర్జీ}} (\emph{alerjī}) is a loan from English \emph{allergy}.
\end{cousinweb}

\begin{grammarlens}[title={What you've built}]
One more word for describing what you lived and felt.
\end{grammarlens}

\subsection*{Your turn}
\begin{itemize}
\raggedright
  \item \emph{Say it:} \emph{alerjī}
  \item \emph{Say it:} \emph{alerjī}, once more
  \item \emph{Say it:} \emph{nāku dummu aṇṭē alerjī.} — I am allergic to dust.
  \item \emph{Recall:} say the Telugu for hearing, then the Telugu for deafness, then say \emph{alerjī} again
\end{itemize}

\begin{tcolorbox}[breakable,colback=teal!4,colframe=teal!35!black,title={Before you move on}]
What does \te{అలెర్జీ} mean? (\textquotedblleft{}An allergy\textquotedblright{}.) Say it once more.
\end{tcolorbox}
\section[\texorpdfstring{\te{నిద్రలేమి} (nidralēmi)}{నిద్రలేమి (nidralēmi)}]{\te{నిద్రలేమి} (nidralēmi) — sleeplessness; insomnia}

\label{lesson:TE-C516-nidralemi}



\begin{quote}
Before the new one: say the Telugu for deafness, then the Telugu for an allergy.
\end{quote}

\subsection*{You'll want to know: \te{నిద్రలేమి}}
\textbf{\te{నిద్రలేమి}} — \emph{nidralēmi} — \textquotedblleft{}sleeplessness; insomnia\textquotedblright{}.

\begin{quote}
\te{పరీక్షల} \te{ముందు} \te{చాలామందికి} \te{నిద్రలేమి} \te{వస్తుంది}.
\end{quote}

\begin{quote}
\emph{parīkṣala mundu cālāmandiki nidralēmi vastundi.}
\end{quote}

\begin{quote}
Before exams many people cannot sleep.
\end{quote}

\te{నిద్ర}, sleep, + \te{లేమి}, lack.

\begin{cousinweb}
\textbf{\te{నిద్రలేమి}} (\emph{nidralēmi}) is built from \te{నిద్ర}, from Sanskrit \emph{nidrā}, \textquotedblleft{}sleep\textquotedblright{} (CDIAL 7200) and the Telugu noun \te{లేమి}, \textquotedblleft{}lack\textquotedblright{}, the noun of \te{లేదు}.
\end{cousinweb}

\begin{grammarlens}[title={What you've built}]
One more word for describing what you lived and felt.
\end{grammarlens}

\subsection*{Your turn}
\begin{itemize}
\raggedright
  \item \emph{Say it:} \emph{nidralēmi}
  \item \emph{Say it:} \emph{nidralēmi}, once more
  \item \emph{Say it:} \emph{parīkṣala mundu cālāmandiki nidralēmi vastundi.} — Before exams many people cannot sleep.
  \item \emph{Recall:} say the Telugu for deafness, then the Telugu for an allergy, then say \emph{nidralēmi} again
\end{itemize}

\begin{tcolorbox}[breakable,colback=teal!4,colframe=teal!35!black,title={Before you move on}]
What does \te{నిద్రలేమి} mean? (\textquotedblleft{}Sleeplessness; insomnia\textquotedblright{}.) Say it once more.
\end{tcolorbox}
